%% Generated by Sphinx.
\def\sphinxdocclass{report}
\documentclass[letterpaper,10pt,english]{sphinxmanual}
\ifdefined\pdfpxdimen
   \let\sphinxpxdimen\pdfpxdimen\else\newdimen\sphinxpxdimen
\fi \sphinxpxdimen=.75bp\relax
\ifdefined\pdfimageresolution
    \pdfimageresolution= \numexpr \dimexpr1in\relax/\sphinxpxdimen\relax
\fi
\newdimen\sphinxremdimen\sphinxremdimen = 10pt
%% let collapsible pdf bookmarks panel have high depth per default
\PassOptionsToPackage{bookmarksdepth=5}{hyperref}

\PassOptionsToPackage{booktabs}{sphinx}
\PassOptionsToPackage{colorrows}{sphinx}

\PassOptionsToPackage{warn}{textcomp}
\usepackage[utf8]{inputenc}
\ifdefined\DeclareUnicodeCharacter
% support both utf8 and utf8x syntaxes
  \ifdefined\DeclareUnicodeCharacterAsOptional
    \def\sphinxDUC#1{\DeclareUnicodeCharacter{"#1}}
  \else
    \let\sphinxDUC\DeclareUnicodeCharacter
  \fi
  \sphinxDUC{00A0}{\nobreakspace}
  \sphinxDUC{2500}{\sphinxunichar{2500}}
  \sphinxDUC{2502}{\sphinxunichar{2502}}
  \sphinxDUC{2514}{\sphinxunichar{2514}}
  \sphinxDUC{251C}{\sphinxunichar{251C}}
  \sphinxDUC{2572}{\textbackslash}
\fi
\usepackage{cmap}
\usepackage[T1]{fontenc}
\usepackage{amsmath,amssymb,amstext}
\usepackage{babel}



\usepackage{tgtermes}
\usepackage{tgheros}
\renewcommand{\ttdefault}{txtt}



\usepackage[Bjarne]{fncychap}
\usepackage{sphinx}

\fvset{fontsize=auto}
\usepackage{geometry}


% Include hyperref last.
\usepackage{hyperref}
% Fix anchor placement for figures with captions.
\usepackage{hypcap}% it must be loaded after hyperref.
% Set up styles of URL: it should be placed after hyperref.
\urlstyle{same}

\addto\captionsenglish{\renewcommand{\contentsname}{Contents:}}

\usepackage{sphinxmessages}
\setcounter{tocdepth}{1}



\title{Software Development Oriented to Machine Learning Practice}
\date{Sep 17, 2026}
\release{}
\author{Ruth Altamirano, Malena Chacón, Odei Martinez de Morentin}
\newcommand{\sphinxlogo}{\vbox{}}
\renewcommand{\releasename}{}
\makeindex
\begin{document}

\ifdefined\shorthandoff
  \ifnum\catcode`\=\string=\active\shorthandoff{=}\fi
  \ifnum\catcode`\"=\active\shorthandoff{"}\fi
\fi

\pagestyle{empty}
\sphinxmaketitle
\pagestyle{plain}
\sphinxtableofcontents
\pagestyle{normal}
\phantomsection\label{\detokenize{index::doc}}


\sphinxAtStartPar
This project focuses on the analysis of asteroid impact risk using machine
learning techniques.

\sphinxAtStartPar
The project includes:
\begin{itemize}
\item {} 
\sphinxAtStartPar
Dataset downloading and preparation.

\item {} 
\sphinxAtStartPar
Data cleaning and transformation.

\item {} 
\sphinxAtStartPar
Feature creation.

\item {} 
\sphinxAtStartPar
Exploratory data analysis.

\item {} 
\sphinxAtStartPar
Neural network training.

\item {} 
\sphinxAtStartPar
Model evaluation.

\item {} 
\sphinxAtStartPar
Impact probability predictions.

\item {} 
\sphinxAtStartPar
Visualization and reporting.

\end{itemize}


\chapter{Project workflow}
\label{\detokenize{index:project-workflow}}
\sphinxAtStartPar
The complete project follows this workflow:

\begin{sphinxVerbatim}[commandchars=\\\{\}]
Dataset from Hugging Face
           |
           v
Dataset downloading and cleaning
           |
           v
Feature creation
           |
           v
Neural network training
           |
           v
Model evaluation
           |
           v
Impact probability predictions
\end{sphinxVerbatim}


\chapter{Project outputs}
\label{\detokenize{index:project-outputs}}
\sphinxAtStartPar
The trained models are stored in \sphinxcode{\sphinxupquote{models/}}.

\sphinxAtStartPar
The metrics, training history, and visualizations are stored in \sphinxcode{\sphinxupquote{reports/}}.

\sphinxAtStartPar
To install and run the project, see {\hyperref[\detokenize{installation::doc}]{\sphinxcrossref{\DUrole{doc}{Installation}}}}.

\sphinxstepscope


\section{Installation}
\label{\detokenize{installation:installation}}\label{\detokenize{installation::doc}}

\subsection{Requirements}
\label{\detokenize{installation:requirements}}
\sphinxAtStartPar
The project requires:
\begin{itemize}
\item {} 
\sphinxAtStartPar
Python 3.10 or higher.

\item {} 
\sphinxAtStartPar
Git.

\item {} 
\sphinxAtStartPar
TensorFlow.

\item {} 
\sphinxAtStartPar
Pandas.

\item {} 
\sphinxAtStartPar
NumPy.

\item {} 
\sphinxAtStartPar
Scikit\sphinxhyphen{}learn.

\item {} 
\sphinxAtStartPar
Joblib.

\item {} 
\sphinxAtStartPar
Typer.

\item {} 
\sphinxAtStartPar
Loguru.

\item {} 
\sphinxAtStartPar
Sphinx, if the documentation needs to be built locally.

\end{itemize}


\subsection{Clone the repository}
\label{\detokenize{installation:clone-the-repository}}
\begin{sphinxVerbatim}[commandchars=\\\{\}]
git\PYG{+w}{ }clone\PYG{+w}{ }https://github.com/altamirano146080/software\PYGZus{}development\PYGZus{}ml\PYGZus{}practice\PYGZus{}1.git
\PYG{n+nb}{cd}\PYG{+w}{ }software\PYGZus{}development\PYGZus{}ml\PYGZus{}practice\PYGZus{}1
\end{sphinxVerbatim}


\subsection{Using \sphinxstyleliteralintitle{\sphinxupquote{uv}}}
\label{\detokenize{installation:using-uv}}
\sphinxAtStartPar
The repository includes \sphinxcode{\sphinxupquote{pyproject.toml}} and \sphinxcode{\sphinxupquote{uv.lock}}.

\sphinxAtStartPar
If \sphinxcode{\sphinxupquote{uv}} is installed, synchronize the project dependencies with:

\begin{sphinxVerbatim}[commandchars=\\\{\}]
uv\PYG{+w}{ }sync
\end{sphinxVerbatim}

\sphinxAtStartPar
Activate the virtual environment if necessary:

\begin{sphinxVerbatim}[commandchars=\\\{\}]
\PYG{n+nb}{source}\PYG{+w}{ }.venv/bin/activate
\end{sphinxVerbatim}


\subsection{Using \sphinxstyleliteralintitle{\sphinxupquote{venv}}}
\label{\detokenize{installation:using-venv}}
\sphinxAtStartPar
Create a virtual environment:

\begin{sphinxVerbatim}[commandchars=\\\{\}]
python\PYG{+w}{ }\PYGZhy{}m\PYG{+w}{ }venv\PYG{+w}{ }.venv
\end{sphinxVerbatim}

\sphinxAtStartPar
On Linux or macOS, activate it with:

\begin{sphinxVerbatim}[commandchars=\\\{\}]
\PYG{n+nb}{source}\PYG{+w}{ }.venv/bin/activate
\end{sphinxVerbatim}

\sphinxAtStartPar
On Windows, activate it with:

\begin{sphinxVerbatim}[commandchars=\\\{\}]
\PYG{p}{.}\PYG{n}{venv}\PYG{p}{\PYGZbs{}}\PYG{n}{Scripts}\PYG{p}{\PYGZbs{}}\PYG{n}{activate}
\end{sphinxVerbatim}

\sphinxAtStartPar
Install the project:

\begin{sphinxVerbatim}[commandchars=\\\{\}]
pip\PYG{+w}{ }install\PYG{+w}{ }\PYGZhy{}e\PYG{+w}{ }.
\end{sphinxVerbatim}


\subsection{Building the documentation}
\label{\detokenize{installation:building-the-documentation}}
\sphinxAtStartPar
Move to the documentation directory:

\begin{sphinxVerbatim}[commandchars=\\\{\}]
\PYG{n+nb}{cd}\PYG{+w}{ }docs
\end{sphinxVerbatim}

\sphinxAtStartPar
Build the HTML documentation:

\begin{sphinxVerbatim}[commandchars=\\\{\}]
make\PYG{+w}{ }html
\end{sphinxVerbatim}

\sphinxAtStartPar
The generated documentation will be available at:

\begin{sphinxVerbatim}[commandchars=\\\{\}]
docs/\PYGZus{}build/html/index.html
\end{sphinxVerbatim}

\sphinxstepscope


\section{Usage}
\label{\detokenize{usage:usage}}\label{\detokenize{usage::doc}}
\sphinxAtStartPar
The project should be executed in the following order:
\begin{enumerate}
\sphinxsetlistlabels{\arabic}{enumi}{enumii}{}{.}%
\item {} 
\sphinxAtStartPar
Download and prepare the dataset.

\item {} 
\sphinxAtStartPar
Create the features and labels.

\item {} 
\sphinxAtStartPar
Train and evaluate the model.

\item {} 
\sphinxAtStartPar
Generate predictions.

\end{enumerate}


\subsection{Download and prepare the dataset}
\label{\detokenize{usage:download-and-prepare-the-dataset}}
\sphinxAtStartPar
The \sphinxcode{\sphinxupquote{dataset}} module downloads the dataset from Hugging Face, creates a
local sample, removes incomplete rows, and saves the cleaned dataset.

\sphinxAtStartPar
Run:

\begin{sphinxVerbatim}[commandchars=\\\{\}]
python\PYG{+w}{ }\PYGZhy{}m\PYG{+w}{ }software\PYGZus{}development\PYGZus{}ml\PYGZus{}practice\PYGZus{}1.dataset
\end{sphinxVerbatim}

\sphinxAtStartPar
By default, the project downloads a sample of up to 500 rows.

\sphinxAtStartPar
To download the dataset again, use:

\begin{sphinxVerbatim}[commandchars=\\\{\}]
python\PYG{+w}{ }\PYGZhy{}m\PYG{+w}{ }software\PYGZus{}development\PYGZus{}ml\PYGZus{}practice\PYGZus{}1.dataset\PYG{+w}{ }\PYGZhy{}\PYGZhy{}force\PYGZhy{}download
\end{sphinxVerbatim}

\sphinxAtStartPar
The generated files are:
\begin{itemize}
\item {} 
\sphinxAtStartPar
\sphinxcode{\sphinxupquote{data/raw/dataset.csv}}

\item {} 
\sphinxAtStartPar
\sphinxcode{\sphinxupquote{data/processed/dataset.csv}}

\end{itemize}


\subsection{Create features}
\label{\detokenize{usage:create-features}}
\sphinxAtStartPar
The \sphinxcode{\sphinxupquote{features}} module creates the input features and the target labels.

\sphinxAtStartPar
Run:

\begin{sphinxVerbatim}[commandchars=\\\{\}]
python\PYG{+w}{ }\PYGZhy{}m\PYG{+w}{ }software\PYGZus{}development\PYGZus{}ml\PYGZus{}practice\PYGZus{}1.features
\end{sphinxVerbatim}

\sphinxAtStartPar
This command creates:
\begin{itemize}
\item {} 
\sphinxAtStartPar
\sphinxcode{\sphinxupquote{data/processed/features.csv}}

\item {} 
\sphinxAtStartPar
\sphinxcode{\sphinxupquote{data/processed/labels.csv}}

\end{itemize}


\subsection{Train the model}
\label{\detokenize{usage:train-the-model}}
\sphinxAtStartPar
The \sphinxcode{\sphinxupquote{train}} module trains and evaluates the neural network.

\sphinxAtStartPar
Run:

\begin{sphinxVerbatim}[commandchars=\\\{\}]
python\PYG{+w}{ }\PYGZhy{}m\PYG{+w}{ }software\PYGZus{}development\PYGZus{}ml\PYGZus{}practice\PYGZus{}1.modeling.train
\end{sphinxVerbatim}

\sphinxAtStartPar
This command creates:
\begin{itemize}
\item {} 
\sphinxAtStartPar
\sphinxcode{\sphinxupquote{models/impact\_probability\_model.keras}}

\item {} 
\sphinxAtStartPar
\sphinxcode{\sphinxupquote{models/feature\_scaler.joblib}}

\item {} 
\sphinxAtStartPar
\sphinxcode{\sphinxupquote{reports/model\_metrics.csv}}

\item {} 
\sphinxAtStartPar
\sphinxcode{\sphinxupquote{reports/training\_history.csv}}

\end{itemize}


\subsection{Generate predictions}
\label{\detokenize{usage:generate-predictions}}
\sphinxAtStartPar
The \sphinxcode{\sphinxupquote{predict}} module loads the trained model and generates predictions.

\sphinxAtStartPar
Run:

\begin{sphinxVerbatim}[commandchars=\\\{\}]
python\PYG{+w}{ }\PYGZhy{}m\PYG{+w}{ }software\PYGZus{}development\PYGZus{}ml\PYGZus{}practice\PYGZus{}1.modeling.predict
\end{sphinxVerbatim}

\sphinxAtStartPar
The predictions are saved as a CSV file in the processed data directory.

\sphinxstepscope


\section{Project Workflow}
\label{\detokenize{workflow:project-workflow}}\label{\detokenize{workflow::doc}}

\subsection{Dataset download}
\label{\detokenize{workflow:dataset-download}}
\sphinxAtStartPar
The project downloads the dataset from Hugging Face using the following
source:

\begin{sphinxVerbatim}[commandchars=\\\{\}]
\PYG{n}{DATASET\PYGZus{}URI} \PYG{o}{=} \PYG{p}{(}
    \PYG{l+s+s2}{\PYGZdq{}}\PYG{l+s+s2}{hf://datasets/juliensimon/sentry\PYGZhy{}impact\PYGZhy{}risk/}\PYG{l+s+s2}{\PYGZdq{}}
    \PYG{l+s+s2}{\PYGZdq{}}\PYG{l+s+s2}{data/sentry\PYGZus{}impact\PYGZus{}risk.parquet}\PYG{l+s+s2}{\PYGZdq{}}
\PYG{p}{)}
\end{sphinxVerbatim}

\sphinxAtStartPar
A sample of up to 500 rows is selected and stored in:

\begin{sphinxVerbatim}[commandchars=\\\{\}]
data/raw/dataset.csv
\end{sphinxVerbatim}

\sphinxAtStartPar
The dataset is then cleaned by removing rows containing missing values.
The cleaned dataset is stored in:

\begin{sphinxVerbatim}[commandchars=\\\{\}]
data/processed/dataset.csv
\end{sphinxVerbatim}


\subsection{Feature engineering}
\label{\detokenize{workflow:feature-engineering}}
\sphinxAtStartPar
The \sphinxcode{\sphinxupquote{features.py}} module reads the processed dataset and removes missing
values if necessary.

\sphinxAtStartPar
The target variable is transformed using a base\sphinxhyphen{}10 logarithm:
\begin{equation*}
\begin{split}y = \log_{10}\left(\text{impact\_probability}\right)\end{split}
\end{equation*}
\sphinxAtStartPar
All numerical columns are selected as model features except:
\begin{itemize}
\item {} 
\sphinxAtStartPar
\sphinxcode{\sphinxupquote{impact\_probability}}

\item {} 
\sphinxAtStartPar
\sphinxcode{\sphinxupquote{log\_impact\_probability}}

\end{itemize}

\sphinxAtStartPar
The resulting files are:
\begin{itemize}
\item {} 
\sphinxAtStartPar
\sphinxcode{\sphinxupquote{features.csv}}: model input variables.

\item {} 
\sphinxAtStartPar
\sphinxcode{\sphinxupquote{labels.csv}}: transformed target variable.

\end{itemize}


\subsection{Model architecture}
\label{\detokenize{workflow:model-architecture}}
\sphinxAtStartPar
The neural network is defined in \sphinxcode{\sphinxupquote{modeling/train.py}}.

\sphinxAtStartPar
It contains:
\begin{itemize}
\item {} 
\sphinxAtStartPar
An input layer.

\item {} 
\sphinxAtStartPar
A dense layer with 64 neurons and ReLU activation.

\item {} 
\sphinxAtStartPar
A dense layer with 32 neurons and ReLU activation.

\item {} 
\sphinxAtStartPar
A dense layer with 16 neurons and ReLU activation.

\item {} 
\sphinxAtStartPar
An output layer with one neuron.

\end{itemize}

\sphinxAtStartPar
The model uses:
\begin{itemize}
\item {} 
\sphinxAtStartPar
The Adam optimizer.

\item {} 
\sphinxAtStartPar
A learning rate of \sphinxcode{\sphinxupquote{0.001}}.

\item {} 
\sphinxAtStartPar
Mean squared error as the loss function.

\item {} 
\sphinxAtStartPar
Mean absolute error as a metric.

\end{itemize}


\subsection{Data splitting and scaling}
\label{\detokenize{workflow:data-splitting-and-scaling}}
\sphinxAtStartPar
The dataset is divided into training and test data using an 80/20 split.

\sphinxAtStartPar
The features are standardized with \sphinxcode{\sphinxupquote{StandardScaler}} before training:

\begin{sphinxVerbatim}[commandchars=\\\{\}]
\PYG{n}{scaler} \PYG{o}{=} \PYG{n}{StandardScaler}\PYG{p}{(}\PYG{p}{)}
\end{sphinxVerbatim}

\sphinxAtStartPar
The scaler is saved to:

\begin{sphinxVerbatim}[commandchars=\\\{\}]
models/feature\PYGZus{}scaler.joblib
\end{sphinxVerbatim}


\subsection{Training}
\label{\detokenize{workflow:training}}
\sphinxAtStartPar
The model is trained for a maximum of 100 epochs with a batch size of 32.

\sphinxAtStartPar
Early stopping is used to reduce overfitting. Training stops when the
validation loss does not improve for 15 consecutive epochs.


\subsection{Evaluation}
\label{\detokenize{workflow:evaluation}}
\sphinxAtStartPar
The model is evaluated using:
\begin{itemize}
\item {} 
\sphinxAtStartPar
MAE: Mean Absolute Error.

\item {} 
\sphinxAtStartPar
RMSE: Root Mean Squared Error.

\item {} 
\sphinxAtStartPar
R\(\sp{\text{2}}\): Coefficient of Determination.

\end{itemize}

\sphinxAtStartPar
The results are saved to:

\begin{sphinxVerbatim}[commandchars=\\\{\}]
reports/model\PYGZus{}metrics.csv
\end{sphinxVerbatim}

\sphinxAtStartPar
The training history is saved to:

\begin{sphinxVerbatim}[commandchars=\\\{\}]
reports/training\PYGZus{}history.csv
\end{sphinxVerbatim}


\subsection{Prediction}
\label{\detokenize{workflow:prediction}}
\sphinxAtStartPar
The prediction module loads:
\begin{itemize}
\item {} 
\sphinxAtStartPar
The trained Keras model.

\item {} 
\sphinxAtStartPar
The saved feature scaler.

\item {} 
\sphinxAtStartPar
The processed feature data.

\end{itemize}

\sphinxAtStartPar
The model produces a prediction in logarithmic form. The original scale is
then recovered using:
\begin{equation*}
\begin{split}\text{impact probability} = 10^{\text{log prediction}}\end{split}
\end{equation*}
\sphinxstepscope


\section{Results}
\label{\detokenize{results:results}}\label{\detokenize{results::doc}}

\subsection{Model metrics}
\label{\detokenize{results:model-metrics}}
\sphinxAtStartPar
The model is evaluated using the following metrics:
\begin{itemize}
\item {} 
\sphinxAtStartPar
Mean Absolute Error.

\item {} 
\sphinxAtStartPar
Root Mean Squared Error.

\item {} 
\sphinxAtStartPar
R\(\sp{\text{2}}\) score.

\end{itemize}

\sphinxAtStartPar
The metrics are stored in:

\begin{sphinxVerbatim}[commandchars=\\\{\}]
reports/model\PYGZus{}metrics.csv
\end{sphinxVerbatim}


\subsection{Training history}
\label{\detokenize{results:training-history}}
\sphinxAtStartPar
The training history is stored in:

\begin{sphinxVerbatim}[commandchars=\\\{\}]
reports/training\PYGZus{}history.csv
\end{sphinxVerbatim}

\sphinxAtStartPar
This file contains information about the training and validation loss during
the training process.


\subsection{Generated visualizations}
\label{\detokenize{results:generated-visualizations}}
\sphinxAtStartPar
The project generates several visualizations in
\sphinxcode{\sphinxupquote{reports/figures/}}.


\subsubsection{Impact probability distribution}
\label{\detokenize{results:impact-probability-distribution}}
\noindent\sphinxincludegraphics[width=700\sphinxpxdimen]{{impact_probability_distribution}.png}


\subsubsection{Training history}
\label{\detokenize{results:id1}}
\noindent\sphinxincludegraphics[width=700\sphinxpxdimen]{{training_history}.png}


\subsubsection{Predictions versus observations}
\label{\detokenize{results:predictions-versus-observations}}
\noindent\sphinxincludegraphics[width=700\sphinxpxdimen]{{predictions_vs_observations}.png}


\subsubsection{Correlation heatmap}
\label{\detokenize{results:correlation-heatmap}}
\noindent\sphinxincludegraphics[width=700\sphinxpxdimen]{{correlation_heatmap}.png}


\subsubsection{Potential impact timeline}
\label{\detokenize{results:potential-impact-timeline}}
\noindent\sphinxincludegraphics[width=700\sphinxpxdimen]{{potential_impact_timeline}.png}


\subsubsection{Magnitude versus Palermo scale}
\label{\detokenize{results:magnitude-versus-palermo-scale}}
\noindent\sphinxincludegraphics[width=700\sphinxpxdimen]{{magnitude_vs_palermo}.png}


\subsubsection{Velocity versus Palermo scale}
\label{\detokenize{results:velocity-versus-palermo-scale}}
\noindent\sphinxincludegraphics[width=700\sphinxpxdimen]{{velocity_vs_palermo}.png}

\sphinxstepscope


\section{Module Reference}
\label{\detokenize{modules:module-reference}}\label{\detokenize{modules::doc}}

\subsection{Dataset}
\label{\detokenize{modules:module-asteroid_impact.dataset}}\label{\detokenize{modules:dataset}}\index{module@\spxentry{module}!software\_development\_ml\_practice\_1.dataset@\spxentry{software\_development\_ml\_practice\_1.dataset}}\index{software\_development\_ml\_practice\_1.dataset@\spxentry{software\_development\_ml\_practice\_1.dataset}!module@\spxentry{module}}\index{download\_dataset() (in module software\_development\_ml\_practice\_1.dataset)@\spxentry{download\_dataset()}\spxextra{in module software\_development\_ml\_practice\_1.dataset}}

\begin{fulllineitems}
\phantomsection\label{\detokenize{modules:asteroid_impact.dataset.download_dataset}}
\pysigstartsignatures
\pysiglinewithargsret
{\sphinxcode{\sphinxupquote{software\_development\_ml\_practice\_1.dataset.}}\sphinxbfcode{\sphinxupquote{download\_dataset}}}
{\sphinxparam{\DUrole{n}{output\_path}\DUrole{p}{:}\DUrole{w}{ }\DUrole{n}{Path}\DUrole{w}{ }\DUrole{o}{=}\DUrole{w}{ }\DUrole{default_value}{PosixPath(\textquotesingle{}/home/ruthaltmno/SoftwareDevelopmentML/software\_development\_ml\_practice\_1/data/raw/dataset.csv\textquotesingle{})}}\sphinxparamcomma \sphinxparam{\DUrole{n}{sample\_size}\DUrole{p}{:}\DUrole{w}{ }\DUrole{n}{int}\DUrole{w}{ }\DUrole{o}{=}\DUrole{w}{ }\DUrole{default_value}{500}}\sphinxparamcomma \sphinxparam{\DUrole{n}{force\_download}\DUrole{p}{:}\DUrole{w}{ }\DUrole{n}{bool}\DUrole{w}{ }\DUrole{o}{=}\DUrole{w}{ }\DUrole{default_value}{False}}}
{{ $\rightarrow$ DataFrame}}
\pysigstopsignatures
\sphinxAtStartPar
Download a sample of the asteroid impact\sphinxhyphen{}risk dataset.

\sphinxAtStartPar
The dataset is downloaded from Hugging Face and stored locally as a CSV
file. If the file already exists, the local copy is reused unless
\sphinxcode{\sphinxupquote{force\_download}} is set to \sphinxcode{\sphinxupquote{True}}.


\subsubsection{Parameters}
\label{\detokenize{modules:parameters}}\begin{description}
\sphinxlineitem{output\_path:}
\sphinxAtStartPar
Path where the downloaded dataset will be saved.

\sphinxlineitem{sample\_size:}
\sphinxAtStartPar
Maximum number of rows to keep in the local sample.

\sphinxlineitem{force\_download:}
\sphinxAtStartPar
Whether to download the dataset again if a local copy exists.

\end{description}


\subsubsection{Returns}
\label{\detokenize{modules:returns}}\begin{description}
\sphinxlineitem{pandas.DataFrame}
\sphinxAtStartPar
The downloaded dataset.

\end{description}

\end{fulllineitems}

\index{main() (in module software\_development\_ml\_practice\_1.dataset)@\spxentry{main()}\spxextra{in module software\_development\_ml\_practice\_1.dataset}}

\begin{fulllineitems}
\phantomsection\label{\detokenize{modules:asteroid_impact.dataset.main}}
\pysigstartsignatures
\pysiglinewithargsret
{\sphinxcode{\sphinxupquote{software\_development\_ml\_practice\_1.dataset.}}\sphinxbfcode{\sphinxupquote{main}}}
{\sphinxparam{\DUrole{n}{force\_download}\DUrole{p}{:}\DUrole{w}{ }\DUrole{n}{bool}\DUrole{w}{ }\DUrole{o}{=}\DUrole{w}{ }\DUrole{default_value}{\textless{}typer.models.OptionInfo object\textgreater{}}}}
{}
\pysigstopsignatures
\sphinxAtStartPar
Downloads and prepares the asteroid impact\sphinxhyphen{}risk dataset.

\end{fulllineitems}

\index{prepare\_dataset() (in module software\_development\_ml\_practice\_1.dataset)@\spxentry{prepare\_dataset()}\spxextra{in module software\_development\_ml\_practice\_1.dataset}}

\begin{fulllineitems}
\phantomsection\label{\detokenize{modules:asteroid_impact.dataset.prepare_dataset}}
\pysigstartsignatures
\pysiglinewithargsret
{\sphinxcode{\sphinxupquote{software\_development\_ml\_practice\_1.dataset.}}\sphinxbfcode{\sphinxupquote{prepare\_dataset}}}
{\sphinxparam{\DUrole{n}{input\_path}\DUrole{p}{:}\DUrole{w}{ }\DUrole{n}{Path}\DUrole{w}{ }\DUrole{o}{=}\DUrole{w}{ }\DUrole{default_value}{PosixPath(\textquotesingle{}/home/ruthaltmno/SoftwareDevelopmentML/software\_development\_ml\_practice\_1/data/raw/dataset.csv\textquotesingle{})}}\sphinxparamcomma \sphinxparam{\DUrole{n}{output\_path}\DUrole{p}{:}\DUrole{w}{ }\DUrole{n}{Path}\DUrole{w}{ }\DUrole{o}{=}\DUrole{w}{ }\DUrole{default_value}{PosixPath(\textquotesingle{}/home/ruthaltmno/SoftwareDevelopmentML/software\_development\_ml\_practice\_1/data/processed/dataset.csv\textquotesingle{})}}}
{{ $\rightarrow$ DataFrame}}
\pysigstopsignatures
\sphinxAtStartPar
Clean the raw dataset and save the processed version.

\sphinxAtStartPar
Rows containing missing values are removed before the processed dataset
is written to disk.


\subsubsection{Parameters}
\label{\detokenize{modules:id1}}\begin{description}
\sphinxlineitem{input\_path:}
\sphinxAtStartPar
Path to the raw dataset.

\sphinxlineitem{output\_path:}
\sphinxAtStartPar
Path where the cleaned dataset will be saved.

\end{description}


\subsubsection{Returns}
\label{\detokenize{modules:id2}}\begin{description}
\sphinxlineitem{pandas.DataFrame}
\sphinxAtStartPar
The cleaned dataset.

\end{description}

\end{fulllineitems}



\subsection{Features}
\label{\detokenize{modules:module-asteroid_impact.features}}\label{\detokenize{modules:features}}\index{module@\spxentry{module}!software\_development\_ml\_practice\_1.features@\spxentry{software\_development\_ml\_practice\_1.features}}\index{software\_development\_ml\_practice\_1.features@\spxentry{software\_development\_ml\_practice\_1.features}!module@\spxentry{module}}
\sphinxAtStartPar
Dataset acquisition and preprocessing utilities.

\sphinxAtStartPar
This module downloads the source dataset from Hugging Face, stores a local
sample, and cleans the data before it is used in feature engineering and model
training.
\index{download\_dataset() (in module software\_development\_ml\_practice\_1.features)@\spxentry{download\_dataset()}\spxextra{in module software\_development\_ml\_practice\_1.features}}

\begin{fulllineitems}
\phantomsection\label{\detokenize{modules:asteroid_impact.features.download_dataset}}
\pysigstartsignatures
\pysiglinewithargsret
{\sphinxcode{\sphinxupquote{software\_development\_ml\_practice\_1.features.}}\sphinxbfcode{\sphinxupquote{download\_dataset}}}
{\sphinxparam{\DUrole{n}{output\_path}\DUrole{p}{:}\DUrole{w}{ }\DUrole{n}{Path}\DUrole{w}{ }\DUrole{o}{=}\DUrole{w}{ }\DUrole{default_value}{PosixPath(\textquotesingle{}/home/ruthaltmno/SoftwareDevelopmentML/software\_development\_ml\_practice\_1/data/raw/dataset.csv\textquotesingle{})}}\sphinxparamcomma \sphinxparam{\DUrole{n}{sample\_size}\DUrole{p}{:}\DUrole{w}{ }\DUrole{n}{int}\DUrole{w}{ }\DUrole{o}{=}\DUrole{w}{ }\DUrole{default_value}{500}}\sphinxparamcomma \sphinxparam{\DUrole{n}{force\_download}\DUrole{p}{:}\DUrole{w}{ }\DUrole{n}{bool}\DUrole{w}{ }\DUrole{o}{=}\DUrole{w}{ }\DUrole{default_value}{False}}}
{{ $\rightarrow$ DataFrame}}
\pysigstopsignatures
\sphinxAtStartPar
Download a sample of the asteroid impact\sphinxhyphen{}risk dataset.

\sphinxAtStartPar
The dataset is downloaded from Hugging Face and stored locally as a CSV
file. If the file already exists, the local copy is reused unless
\sphinxcode{\sphinxupquote{force\_download}} is set to \sphinxcode{\sphinxupquote{True}}.


\subsubsection{Parameters}
\label{\detokenize{modules:id3}}\begin{description}
\sphinxlineitem{output\_path}{[}Path, default=RAW\_DATA\_PATH{]}
\sphinxAtStartPar
Path where the downloaded dataset will be saved.

\sphinxlineitem{sample\_size}{[}int, default=500{]}
\sphinxAtStartPar
Maximum number of rows to keep in the local sample.

\sphinxlineitem{force\_download}{[}bool, default=False{]}
\sphinxAtStartPar
Whether to download the dataset again if a local copy exists.

\end{description}


\subsubsection{Returns}
\label{\detokenize{modules:id4}}\begin{description}
\sphinxlineitem{pandas.DataFrame}
\sphinxAtStartPar
The downloaded dataset.

\end{description}

\end{fulllineitems}

\index{main() (in module software\_development\_ml\_practice\_1.features)@\spxentry{main()}\spxextra{in module software\_development\_ml\_practice\_1.features}}

\begin{fulllineitems}
\phantomsection\label{\detokenize{modules:asteroid_impact.features.main}}
\pysigstartsignatures
\pysiglinewithargsret
{\sphinxcode{\sphinxupquote{software\_development\_ml\_practice\_1.features.}}\sphinxbfcode{\sphinxupquote{main}}}
{\sphinxparam{\DUrole{n}{force\_download}\DUrole{p}{:}\DUrole{w}{ }\DUrole{n}{bool}\DUrole{w}{ }\DUrole{o}{=}\DUrole{w}{ }\DUrole{default_value}{\textless{}typer.models.OptionInfo object\textgreater{}}}}
{}
\pysigstopsignatures
\sphinxAtStartPar
Download and prepare the asteroid impact\sphinxhyphen{}risk dataset.


\subsubsection{Parameters}
\label{\detokenize{modules:id5}}\begin{description}
\sphinxlineitem{force\_download}{[}bool, default=False{]}
\sphinxAtStartPar
If True, redownload the source dataset even when a local copy already
exists.

\end{description}


\subsubsection{Returns}
\label{\detokenize{modules:id6}}\begin{description}
\sphinxlineitem{None}
\sphinxAtStartPar
This command runs the data acquisition and preprocessing pipeline.

\end{description}

\end{fulllineitems}

\index{prepare\_dataset() (in module software\_development\_ml\_practice\_1.features)@\spxentry{prepare\_dataset()}\spxextra{in module software\_development\_ml\_practice\_1.features}}

\begin{fulllineitems}
\phantomsection\label{\detokenize{modules:asteroid_impact.features.prepare_dataset}}
\pysigstartsignatures
\pysiglinewithargsret
{\sphinxcode{\sphinxupquote{software\_development\_ml\_practice\_1.features.}}\sphinxbfcode{\sphinxupquote{prepare\_dataset}}}
{\sphinxparam{\DUrole{n}{input\_path}\DUrole{p}{:}\DUrole{w}{ }\DUrole{n}{Path}\DUrole{w}{ }\DUrole{o}{=}\DUrole{w}{ }\DUrole{default_value}{PosixPath(\textquotesingle{}/home/ruthaltmno/SoftwareDevelopmentML/software\_development\_ml\_practice\_1/data/raw/dataset.csv\textquotesingle{})}}\sphinxparamcomma \sphinxparam{\DUrole{n}{output\_path}\DUrole{p}{:}\DUrole{w}{ }\DUrole{n}{Path}\DUrole{w}{ }\DUrole{o}{=}\DUrole{w}{ }\DUrole{default_value}{PosixPath(\textquotesingle{}/home/ruthaltmno/SoftwareDevelopmentML/software\_development\_ml\_practice\_1/data/processed/dataset.csv\textquotesingle{})}}}
{{ $\rightarrow$ DataFrame}}
\pysigstopsignatures
\sphinxAtStartPar
Clean the raw dataset and save the processed version.

\sphinxAtStartPar
Rows containing missing values are removed before the processed dataset
is written to disk.


\subsubsection{Parameters}
\label{\detokenize{modules:id7}}\begin{description}
\sphinxlineitem{input\_path}{[}Path, default=RAW\_DATA\_PATH{]}
\sphinxAtStartPar
Path to the raw dataset.

\sphinxlineitem{output\_path}{[}Path, default=PROCESSED\_DATA\_PATH{]}
\sphinxAtStartPar
Path where the cleaned dataset will be saved.

\end{description}


\subsubsection{Returns}
\label{\detokenize{modules:id8}}\begin{description}
\sphinxlineitem{pandas.DataFrame}
\sphinxAtStartPar
The cleaned dataset.

\end{description}

\end{fulllineitems}



\subsection{Configuration}
\label{\detokenize{modules:module-asteroid_impact.config}}\label{\detokenize{modules:configuration}}\index{module@\spxentry{module}!software\_development\_ml\_practice\_1.config@\spxentry{software\_development\_ml\_practice\_1.config}}\index{software\_development\_ml\_practice\_1.config@\spxentry{software\_development\_ml\_practice\_1.config}!module@\spxentry{module}}
\sphinxAtStartPar
Configuration settings for the asteroid impact\sphinxhyphen{}risk project.

\sphinxAtStartPar
This module defines the project root and the locations used for raw data,
processed data, models, and generated reports. It also ensures the required
folders exist when the package is imported.


\subsection{Plots}
\label{\detokenize{modules:module-asteroid_impact.plots}}\label{\detokenize{modules:plots}}\index{module@\spxentry{module}!software\_development\_ml\_practice\_1.plots@\spxentry{software\_development\_ml\_practice\_1.plots}}\index{software\_development\_ml\_practice\_1.plots@\spxentry{software\_development\_ml\_practice\_1.plots}!module@\spxentry{module}}
\sphinxAtStartPar
Feature engineering utilities for the asteroid impact\sphinxhyphen{}risk model.

\sphinxAtStartPar
The module creates the feature matrix and transformed target variable used to
train the regression model.
\index{create\_features() (in module software\_development\_ml\_practice\_1.plots)@\spxentry{create\_features()}\spxextra{in module software\_development\_ml\_practice\_1.plots}}

\begin{fulllineitems}
\phantomsection\label{\detokenize{modules:asteroid_impact.plots.create_features}}
\pysigstartsignatures
\pysiglinewithargsret
{\sphinxcode{\sphinxupquote{software\_development\_ml\_practice\_1.plots.}}\sphinxbfcode{\sphinxupquote{create\_features}}}
{\sphinxparam{\DUrole{n}{input\_path}\DUrole{p}{:}\DUrole{w}{ }\DUrole{n}{Path}\DUrole{w}{ }\DUrole{o}{=}\DUrole{w}{ }\DUrole{default_value}{PosixPath(\textquotesingle{}/home/ruthaltmno/SoftwareDevelopmentML/software\_development\_ml\_practice\_1/data/processed/dataset.csv\textquotesingle{})}}\sphinxparamcomma \sphinxparam{\DUrole{n}{features\_path}\DUrole{p}{:}\DUrole{w}{ }\DUrole{n}{Path}\DUrole{w}{ }\DUrole{o}{=}\DUrole{w}{ }\DUrole{default_value}{PosixPath(\textquotesingle{}/home/ruthaltmno/SoftwareDevelopmentML/software\_development\_ml\_practice\_1/data/processed/features.csv\textquotesingle{})}}\sphinxparamcomma \sphinxparam{\DUrole{n}{labels\_path}\DUrole{p}{:}\DUrole{w}{ }\DUrole{n}{Path}\DUrole{w}{ }\DUrole{o}{=}\DUrole{w}{ }\DUrole{default_value}{PosixPath(\textquotesingle{}/home/ruthaltmno/SoftwareDevelopmentML/software\_development\_ml\_practice\_1/data/processed/labels.csv\textquotesingle{})}}}
{{ $\rightarrow$ tuple\DUrole{p}{{[}}DataFrame\DUrole{p}{,}\DUrole{w}{ }Series\DUrole{p}{{]}}}}
\pysigstopsignatures
\sphinxAtStartPar
Create model features and the transformed target variable.

\sphinxAtStartPar
Missing rows are removed and the impact probability is transformed using
a base\sphinxhyphen{}10 logarithm. Numerical columns are used as model features.


\subsubsection{Parameters}
\label{\detokenize{modules:id9}}\begin{description}
\sphinxlineitem{input\_path}{[}Path, default=DATASET\_PATH{]}
\sphinxAtStartPar
Path to the processed dataset.

\sphinxlineitem{features\_path}{[}Path, default=FEATURES\_PATH{]}
\sphinxAtStartPar
Path where the feature data will be saved.

\sphinxlineitem{labels\_path}{[}Path, default=LABELS\_PATH{]}
\sphinxAtStartPar
Path where the target labels will be saved.

\end{description}


\subsubsection{Returns}
\label{\detokenize{modules:id10}}\begin{description}
\sphinxlineitem{tuple{[}pandas.DataFrame, pandas.Series{]}}
\sphinxAtStartPar
The feature matrix and transformed target variable.

\end{description}

\end{fulllineitems}

\index{main() (in module software\_development\_ml\_practice\_1.plots)@\spxentry{main()}\spxextra{in module software\_development\_ml\_practice\_1.plots}}

\begin{fulllineitems}
\phantomsection\label{\detokenize{modules:asteroid_impact.plots.main}}
\pysigstartsignatures
\pysiglinewithargsret
{\sphinxcode{\sphinxupquote{software\_development\_ml\_practice\_1.plots.}}\sphinxbfcode{\sphinxupquote{main}}}
{}
{}
\pysigstopsignatures
\sphinxAtStartPar
Generate the feature and label files.


\subsubsection{Returns}
\label{\detokenize{modules:id11}}\begin{description}
\sphinxlineitem{None}
\sphinxAtStartPar
Executes the feature creation pipeline.

\end{description}

\end{fulllineitems}



\subsection{Model training}
\label{\detokenize{modules:module-asteroid_impact.modeling.train}}\label{\detokenize{modules:model-training}}\index{module@\spxentry{module}!software\_development\_ml\_practice\_1.modeling.train@\spxentry{software\_development\_ml\_practice\_1.modeling.train}}\index{software\_development\_ml\_practice\_1.modeling.train@\spxentry{software\_development\_ml\_practice\_1.modeling.train}!module@\spxentry{module}}
\sphinxAtStartPar
Modeling package for training and inference utilities.


\subsection{Model prediction}
\label{\detokenize{modules:module-asteroid_impact.modeling.predict}}\label{\detokenize{modules:model-prediction}}\index{module@\spxentry{module}!software\_development\_ml\_practice\_1.modeling.predict@\spxentry{software\_development\_ml\_practice\_1.modeling.predict}}\index{software\_development\_ml\_practice\_1.modeling.predict@\spxentry{software\_development\_ml\_practice\_1.modeling.predict}!module@\spxentry{module}}
\sphinxAtStartPar
Model training pipeline for the impact\sphinxhyphen{}probability regressor.

\sphinxAtStartPar
This module builds the neural network, fits it on scaled features, evaluates
performance, and saves the trained model along with training metrics.
\index{build\_model() (in module software\_development\_ml\_practice\_1.modeling.predict)@\spxentry{build\_model()}\spxextra{in module software\_development\_ml\_practice\_1.modeling.predict}}

\begin{fulllineitems}
\phantomsection\label{\detokenize{modules:asteroid_impact.modeling.predict.build_model}}
\pysigstartsignatures
\pysiglinewithargsret
{\sphinxcode{\sphinxupquote{software\_development\_ml\_practice\_1.modeling.predict.}}\sphinxbfcode{\sphinxupquote{build\_model}}}
{\sphinxparam{\DUrole{n}{input\_shape}\DUrole{p}{:}\DUrole{w}{ }\DUrole{n}{int}}}
{{ $\rightarrow$ Model}}
\pysigstopsignatures
\sphinxAtStartPar
Create the neural network architecture.


\subsubsection{Parameters}
\label{\detokenize{modules:id12}}\begin{description}
\sphinxlineitem{input\_shape}{[}int{]}
\sphinxAtStartPar
Number of input features.

\end{description}


\subsubsection{Returns}
\label{\detokenize{modules:id13}}\begin{description}
\sphinxlineitem{tensorflow.keras.Model}
\sphinxAtStartPar
A compiled neural network model.

\end{description}

\end{fulllineitems}

\index{main() (in module software\_development\_ml\_practice\_1.modeling.predict)@\spxentry{main()}\spxextra{in module software\_development\_ml\_practice\_1.modeling.predict}}

\begin{fulllineitems}
\phantomsection\label{\detokenize{modules:asteroid_impact.modeling.predict.main}}
\pysigstartsignatures
\pysiglinewithargsret
{\sphinxcode{\sphinxupquote{software\_development\_ml\_practice\_1.modeling.predict.}}\sphinxbfcode{\sphinxupquote{main}}}
{}
{}
\pysigstopsignatures
\sphinxAtStartPar
Train and evaluate the model.


\subsubsection{Returns}
\label{\detokenize{modules:id14}}\begin{description}
\sphinxlineitem{None}
\sphinxAtStartPar
Executes the training pipeline.

\end{description}

\end{fulllineitems}

\index{train\_model() (in module software\_development\_ml\_practice\_1.modeling.predict)@\spxentry{train\_model()}\spxextra{in module software\_development\_ml\_practice\_1.modeling.predict}}

\begin{fulllineitems}
\phantomsection\label{\detokenize{modules:asteroid_impact.modeling.predict.train_model}}
\pysigstartsignatures
\pysiglinewithargsret
{\sphinxcode{\sphinxupquote{software\_development\_ml\_practice\_1.modeling.predict.}}\sphinxbfcode{\sphinxupquote{train\_model}}}
{\sphinxparam{\DUrole{n}{features\_path}\DUrole{p}{:}\DUrole{w}{ }\DUrole{n}{Path}\DUrole{w}{ }\DUrole{o}{=}\DUrole{w}{ }\DUrole{default_value}{PosixPath(\textquotesingle{}/home/ruthaltmno/SoftwareDevelopmentML/software\_development\_ml\_practice\_1/data/processed/features.csv\textquotesingle{})}}\sphinxparamcomma \sphinxparam{\DUrole{n}{labels\_path}\DUrole{p}{:}\DUrole{w}{ }\DUrole{n}{Path}\DUrole{w}{ }\DUrole{o}{=}\DUrole{w}{ }\DUrole{default_value}{PosixPath(\textquotesingle{}/home/ruthaltmno/SoftwareDevelopmentML/software\_development\_ml\_practice\_1/data/processed/labels.csv\textquotesingle{})}}\sphinxparamcomma \sphinxparam{\DUrole{n}{model\_path}\DUrole{p}{:}\DUrole{w}{ }\DUrole{n}{Path}\DUrole{w}{ }\DUrole{o}{=}\DUrole{w}{ }\DUrole{default_value}{PosixPath(\textquotesingle{}/home/ruthaltmno/SoftwareDevelopmentML/software\_development\_ml\_practice\_1/models/impact\_probability\_model.keras\textquotesingle{})}}\sphinxparamcomma \sphinxparam{\DUrole{n}{scaler\_path}\DUrole{p}{:}\DUrole{w}{ }\DUrole{n}{Path}\DUrole{w}{ }\DUrole{o}{=}\DUrole{w}{ }\DUrole{default_value}{PosixPath(\textquotesingle{}/home/ruthaltmno/SoftwareDevelopmentML/software\_development\_ml\_practice\_1/models/feature\_scaler.joblib\textquotesingle{})}}\sphinxparamcomma \sphinxparam{\DUrole{n}{metrics\_path}\DUrole{p}{:}\DUrole{w}{ }\DUrole{n}{Path}\DUrole{w}{ }\DUrole{o}{=}\DUrole{w}{ }\DUrole{default_value}{PosixPath(\textquotesingle{}/home/ruthaltmno/SoftwareDevelopmentML/software\_development\_ml\_practice\_1/reports/model\_metrics.csv\textquotesingle{})}}}
{{ $\rightarrow$ None}}
\pysigstopsignatures
\sphinxAtStartPar
Train, evaluate, and save the neural network model.

\sphinxAtStartPar
The function splits the data, scales the features, trains the model,
calculates evaluation metrics, and saves the model and scaler.


\subsubsection{Parameters}
\label{\detokenize{modules:id15}}\begin{description}
\sphinxlineitem{features\_path}{[}Path, default=FEATURES\_PATH{]}
\sphinxAtStartPar
Path to the input feature data.

\sphinxlineitem{labels\_path}{[}Path, default=LABELS\_PATH{]}
\sphinxAtStartPar
Path to the target labels.

\sphinxlineitem{model\_path}{[}Path, default=MODEL\_PATH{]}
\sphinxAtStartPar
Path where the trained model will be saved.

\sphinxlineitem{scaler\_path}{[}Path, default=SCALER\_PATH{]}
\sphinxAtStartPar
Path where the feature scaler will be saved.

\sphinxlineitem{metrics\_path}{[}Path, default=METRICS\_PATH{]}
\sphinxAtStartPar
Path where the evaluation metrics will be saved.

\end{description}


\subsubsection{Returns}
\label{\detokenize{modules:id16}}\begin{description}
\sphinxlineitem{None}
\sphinxAtStartPar
Writes the trained model, scaler, and metrics to disk.

\end{description}

\end{fulllineitems}



\renewcommand{\indexname}{Python Module Index}
\begin{sphinxtheindex}
\let\bigletter\sphinxstyleindexlettergroup
\bigletter{s}
\item\relax\sphinxstyleindexentry{software\_development\_ml\_practice\_1.config}\sphinxstyleindexpageref{modules:\detokenize{module-asteroid_impact.config}}
\item\relax\sphinxstyleindexentry{software\_development\_ml\_practice\_1.dataset}\sphinxstyleindexpageref{modules:\detokenize{module-asteroid_impact.dataset}}
\item\relax\sphinxstyleindexentry{software\_development\_ml\_practice\_1.features}\sphinxstyleindexpageref{modules:\detokenize{module-asteroid_impact.features}}
\item\relax\sphinxstyleindexentry{software\_development\_ml\_practice\_1.modeling.predict}\sphinxstyleindexpageref{modules:\detokenize{module-asteroid_impact.modeling.predict}}
\item\relax\sphinxstyleindexentry{software\_development\_ml\_practice\_1.modeling.train}\sphinxstyleindexpageref{modules:\detokenize{module-asteroid_impact.modeling.train}}
\item\relax\sphinxstyleindexentry{software\_development\_ml\_practice\_1.plots}\sphinxstyleindexpageref{modules:\detokenize{module-asteroid_impact.plots}}
\end{sphinxtheindex}

\renewcommand{\indexname}{Index}
\printindex
\end{document}